Table~\ref{tab:repeated-corruption} gives the complete budget/threshold comparison. A useful conditional result occurs on loans at budget 64 and $q=.99$: generation raises mean fully recovered events from 30.3 to 67.0 with zero observed new damage in either arm. It improves recovery in two seeds and loses in one; mean history F1 is slightly lower. At $q=.95$, the same augmentation reduces recovery from 236.7 to 108.3 at equal mean damage. Purchasing at budget 64 trades fewer recoveries for less damage at .95, but loses on both measures at .99. At budget 256, augmentation lowers mean history F1 in both sources at both targets. This fixed mixture therefore offers no established general advantage beyond repeated corruption; neither a universal verification-budget saving nor the optimal point to stop using generation follows from these trials.

Table~\ref{tab:followup-ablation} removes information before retrieval and refitting. On loans, removing timestamps reduces mean full-event recovery from 108.3 to 7.7, also reducing damage. Purchasing loses history F1 without context objects or descriptions. Descriptions are not uniformly beneficial: removing them slightly improves loan F1 and recovery while increasing damage. Removing all history edges yields no repairs. These results identify information dependencies of the scorer and candidate pipeline; they do not establish superiority of a graph neural architecture.

At budget 64 and .99, the guard reduces real-only loan damage from 39.7 to 0.0 events while preserving 30.3 recoveries, but reduces correctly merged identity pairs from 5.3 to 1.7. For augmented purchasing, damage falls from 10.3 to 4.7 with recovery changing from 307.7 to 308.7; some false pairs remain. All 192 empty-history model--event decisions remain unresolved both with and without the guard, so this challenge establishes no incremental guard benefit. Guarded augmented purchasing still damages 5.3 events on clean inputs and recovers only 54.0 events while damaging 39.3 under pipeline faults. Pipeline faults appear in training; this is a known-family stress test.

On a Threadripper PRO 5955WX with four fitting threads per worker, budget-64 real-only/augmented fitting takes mean 3.5/5.7 seconds for loans and 4.6/7.1 for purchasing. The shared generated feature cache costs 10.1 seconds; full-input inference including candidate construction and metric computation takes about 33/21 seconds per loan/purchasing test graph, excluding file writes. These are measured host timings, not human verification savings. All 60 calibration computations and 156 saved repairs were replayed; independent checks verified history metrics, cluster-expanded harm, event preservation and exact undo. Per-seed ranges and full costs remain in the artifact; the artifact includes the fixed recovery--damage frontier.
